%=====================================================================
\section{Scope, hypotheses and notation}
\label{sec:overview}

THUMS stores electrical energy as latent heat in a phase-change material packed
into the annulus of a repurposed oil-and-gas well. A high-temperature heat pump
melts the PCM during charging; an organic Rankine cycle recovers the energy
during discharging. The storage medium is a trimodal PCM melting at
\SI{150}{\celsius}; the heat-transfer fluid is pressurised water; the borehole
contains longitudinally finned hairpin tubes.

The model answers two questions: \emph{how many wells does a given power block
need}, and \emph{what round-trip efficiency does the resulting plant achieve}.
Everything below serves those two numbers.

%---------------------------------------------------------------------
\subsection*{Headline results, version 0.4}

\begin{center}
\begin{tabular}{@{}l r r r@{}}
\toprule
 & v0.1 published & marched, old $k_{w}$ & \textbf{v0.4} \\
\midrule
$\Nwells$                     & 24.123 & 22.431 & \textbf{11.573} \\
binding constraint            & --- & charge rate & \textbf{capacity} \\
$\varepsilon_{\mathrm{PCM}}$  & 0.5200 & 0.5402 & \textbf{0.9992} \\
local $\varepsilon$, min--max & --- & 0.46--1.78 & \textbf{0.922--1.000} \\
$\eta_{\mathrm{RTE}}$         & 0.43027 & 0.42955 & 0.41585 \\
$\eta_{\mathrm{RTE}}^{0}$     & 0.43484 & 0.43484 & 0.43484 \\
$\mathrm{COP}$                & 2.95633 & 2.95633 & 2.95633 \\
$\eta_{\mathrm{ORC}}$         & 0.23685 & 0.23685 & 0.23685 \\
$\eta_{\mathrm{storage}}$     & --- & 0.953 & 0.9364 \\
closure error                 & 3--54\,\% & $10^{-16}$ & $10^{-16}$ \\
\bottomrule
\end{tabular}
\end{center}

\noindent The sizing moves a long way and the efficiency hardly at all:
$\Nwells$ falls by \SI{52}{\percent} while $\eta_{\mathrm{RTE}}$ moves by about
\SI{3}{\percent}. $\eta_{\mathrm{RTE}}^{0}$ is identical to twelve significant
figures throughout --- the cycles are untouched by any of this. What the storage
model determines is how much hardware the cycle needs, not how efficiently it
runs.

%---------------------------------------------------------------------
\subsection{The structure of the model}

The calculation is a chain, not a coupled system. Each stage consumes the
previous stage's output and nothing flows backwards:

\begin{enumerate}[itemsep=1pt]
\item the ORC and heat-pump cycles fix $\eta_{\mathrm{ORC}}$ and
      $\mathrm{COP}$ from CoolProp property calls alone
      (\S\ref{sec:cycles:states});
\item an energy budget converts the specified electrical output into a charging
      heat duty and a secondary-fluid mass flow (\S\ref{sec:cycles:budget});
\item a resistance network gives the local overall coefficient $\Ui$ along the
      tube (\S\ref{sec:network});
\item an effectiveness--NTU march integrates the fluid temperature down the
      tube, segment by segment (\S\ref{sec:front:march});
\item a melt-front model converts the delivered heat into a melted volume
      (\S\ref{sec:front:legacy} or \S\ref{sec:front:marched});
\item the well count is chosen so that the field satisfies its constraints, and
      the performance indices follow (\S\ref{sec:hydraulics}).
\end{enumerate}

The consequence worth stating explicitly is that the cycles are decoupled from
the storage. $\eta_{\mathrm{ORC}}$ and $\mathrm{COP}$ do not depend on anything
the borehole does. No change to the storage model can move them, which is why
the ideal round-trip efficiency $\eta_{\mathrm{RTE}}^{0}$ of
Eq.~\eqref{eq:rte0} is invariant across every formulation in this document.

%---------------------------------------------------------------------
\subsection{Modelling hypotheses}
\label{sec:overview:hypotheses}

The following are assumed throughout. Those known to be violated in the
operating regime are revisited in \S\ref{sec:defects}.

\begin{description}[leftmargin=1.6em,itemsep=2pt]
\item[H1 --- Quasi-steady conduction.] The melt layer carries no thermal mass.
  Its temperature profile is the steady cylindrical profile corresponding to the
  instantaneous front position. Time enters the model \emph{only} through the
  front position. The hypothesis is controlled by the Stefan number
  \begin{equation}
    \mathrm{Ste} \;=\; \frac{c_{p,m}\,\Delta T}{\hm},
    \label{eq:stefan-number}
  \end{equation}
  and degrades as it grows; at the design point $\mathrm{Ste}=0.0474$.
\item[H2 --- No axial conduction.] Neither the tube wall nor the PCM conducts
  along the well. Segments are coupled only through the fluid stream. This one
  is safe by a wide margin: the axial diffusion time over a segment is
  $\Delta z^{2}/\alpha_{l}=\SI{5.4e9}{\second}$ against a
  \SI{3.6e4}{\second} window (\S\ref{sec:front:segeq}).
\item[H3 --- Concentric annular front.] The melt region around each tube is a
  cylindrical annulus of uniform thickness $\delta$, so that the melted
  cross-section is $\pi[(r_e+\delta)^2-r_e^2]$. Fronts from neighbouring tubes
  do not interact and do not merge. \emph{Under Formulation~C this cannot be
  violated, but it is saturated over the whole well --- see
  \S\ref{sec:defects:h3}.}
\item[H4 --- Conduction only in the melt.] No natural convection in the liquid
  PCM. This is conservative early and increasingly wrong as $\delta$ grows.
\item[H5 --- Enthalpy state, no radial resolution.] \emph{Revised in v0.4.} The
  PCM carries latent \emph{and} sensible energy, lumped into one enthalpy per
  segment (\S\ref{sec:front:enthalpy}): fully molten PCM superheats, fully solid
  PCM subcools. What is still absent is any radial resolution of temperature
  \emph{within} a phase --- each segment has a single PCM temperature, so the
  thermal gradient through the melt layer is quasi-steady (H1) rather than
  resolved.
\item[H6 --- Adiabatic borehole wall.] No heat exchange with the formation in
  either direction, and no geothermal gradient in the energy balance.
\item[H7 --- Incompressible, single-phase secondary fluid.] Pressurised water at
  constant $P$; no buoyancy head in the vertical loop. The fluid is also taken
  as \emph{quasi-steady}: $\partial T/\partial t$ is dropped from its energy
  equation, which costs an unresolved startup transient of about
  \SI{38}{\minute} in each half-cycle (\S\ref{sec:front:segeq}).
\item[H8 --- Multilayer PCM.] The well is divided into $N_{\mathrm{lay}}$ axial
  layers of distinct melting temperature, so that the PCM melting point follows
  the fluid glide down the well. Layer boundaries are equally spaced along the
  \emph{developed} tube length $L_{\mathrm{tube}}$, not the depth --- so a
  hairpin's down leg and return leg belong to \emph{different} layers at the
  same depth. See H9.
\item[H9 --- Perfect azimuthal insulation.] \emph{New in v0.4a, and the
  weakest assumption in the set.} Each tube leg owns an equivalent PCM cell of
  radius $r_{\mathrm{cell}}$ (Eq.~\eqref{eq:rcell}) whose outer boundary is
  adiabatic: no heat crosses between the cells of neighbouring legs at the same
  depth. Because H8 cascades along the developed length, those neighbours are at
  \emph{different} melting temperatures --- up to \SI{48.9}{\kelvin} apart at
  the wellhead, falling to zero at the well bottom where the two legs meet. The
  assumption therefore requires a physical partition between the down-leg and
  return-leg PCM \emph{and} that the partition be insulating. Neither is
  designed, costed, or modelled. A slab estimate puts the neglected leak at
  $\sim\SI{22}{\watt\per\meter}$ against a useful duty of
  $\sim\SI{83}{\watt\per\meter}$ at the wellhead --- of order
  \SI{25}{\percent}, and acting to erase precisely the temperature separation
  the cascade exists to create. Held as an idealisation pending evaluation;
  \S\ref{sec:defects:h9}.
\end{description}

%---------------------------------------------------------------------
\subsection{Geometry}
\label{sec:overview:geometry}

One borehole of diameter $D_{\mathrm{well}}$ and depth $L_{\mathrm{well}}$
contains $n_{t}$ hairpin tubes. A hairpin runs down and back, so its developed
length is
\begin{equation}
  L_{\mathrm{tube}} \;=\; 2\,L_{\mathrm{well}},
  \label{eq:Ltube}
\end{equation}
and the number of tube \emph{legs} in the borehole is $2n_{t}$. This
factor-of-two bookkeeping is a recurring source of error: per-borehole
quantities multiply the per-tube result by $n_{t}$, not by $2n_{t}$, because
Eq.~\eqref{eq:Ltube} has already counted both legs.

The distinction between the developed coordinate $s\in[0,L_{\mathrm{tube}}]$
and the depth $d$ matters wherever a profile is plotted or a layer assigned:
\begin{equation}
  d(s) =
  \begin{cases}
    s, & s \le L_{\mathrm{well}} \quad\text{(down leg)},\\
    L_{\mathrm{tube}} - s, & s > L_{\mathrm{well}} \quad\text{(return leg)}.
  \end{cases}
  \label{eq:depth-fold}
\end{equation}
Every depth therefore carries \emph{two} legs of the same hairpin, at developed
positions $s$ and $L_{\mathrm{tube}}-s$, and by H8 those two are in different
PCM layers. H9 assumes they do not exchange heat.

Each leg carries $n_{f}$ longitudinal fins of thickness $t_{f}$ and radial
length $L_{f}$. The PCM volume available in one borehole is the borehole volume
less the tubes and their fins,
\begin{equation}
  V_{\mathrm{well}}
  \;=\;
  \frac{\pi D_{\mathrm{well}}^{2}}{4}\,L_{\mathrm{well}}
  \;-\;
  2n_{t}\Bigl(\pi r_{e}^{2} + n_{f} t_{f} L_{f}\Bigr) L_{\mathrm{well}} .
  \label{eq:Vwell}
\end{equation}

Equation~\eqref{eq:Vwell} deserves a second look, because it is the source of a
reporting trap taken up in \S\ref{sec:hydraulics:capacity}: fin metal displaces
PCM, so \emph{any} change to the fins moves $V_{\mathrm{well}}$ and hence the
storage capacity of a borehole, entirely independently of what the fins do
thermally.

\emph{Source:} \code{thums\_core/config.py}, class \code{Geometry}.

%---------------------------------------------------------------------
\subsection{Nomenclature}

\begin{longtable}{@{}l l l@{}}
\toprule
symbol & meaning & unit \\
\midrule
\endhead
$\Amelt$ & melted cross-section per unit tube length & \si{\meter\squared} \\
$A_{\mathrm{flow}}$ & tube flow area, $\pi r_{i}^{2}$ & \si{\meter\squared} \\
$d$ & depth below surface & \si{\meter} \\
$s$ & developed tube coordinate & \si{\meter} \\
$A_{\mathrm{avail}}$ & PCM cross-section owned by one segment & \si{\meter\squared} \\
$C_{s}$, $C_{l}$ & sensible capacity per unit tube length & \si{\joule\per\meter\per\kelvin} \\
$E'$ & PCM enthalpy per unit tube length & \si{\joule\per\meter} \\
$E'_{\mathrm{lat}}$ & latent capacity per unit tube length & \si{\joule\per\meter} \\
$K$ & segment conductance per unit tube length & \si{\watt\per\meter\per\kelvin} \\
$c_{p}$ & specific heat capacity & \si{\joule\per\kilogram\per\kelvin} \\
$D_{\mathrm{well}}$ & borehole diameter & \si{\meter} \\
$\hm$ & latent heat of fusion & \si{\joule\per\kilogram} \\
$h_{i}$ & internal convective coefficient & \si{\watt\per\meter\squared\per\kelvin} \\
$h_{e}$ & melt-layer conductance, outer side & \si{\watt\per\meter\squared\per\kelvin} \\
$\km$ & PCM thermal conductivity & \si{\watt\per\meter\per\kelvin} \\
$k_{w}$ & tube-wall and fin conductivity & \si{\watt\per\meter\per\kelvin} \\
$L_{f}$, $t_{f}$, $n_{f}$ & fin length, thickness, count & \si{\meter}, \si{\meter}, -- \\
$L_{\mathrm{well}}$, $L_{\mathrm{tube}}$ & well depth, developed tube length & \si{\meter} \\
$\md$ & mass flow rate & \si{\kilogram\per\second} \\
$n_{t}$ & hairpin tubes per borehole & -- \\
$N_{\mathrm{lay}}$ & PCM layers along the well & -- \\
$\Nwells$ & number of boreholes & -- \\
$\mathrm{NTU}$ & number of transfer units, per segment & -- \\
$q'$ & heat rate per unit tube length & \si{\watt\per\meter} \\
$r_{i}$, $r_{e}$ & tube inner and outer radius & \si{\meter} \\
$r_{\mathrm{cell}}$ & equivalent cell radius of one tube leg & \si{\meter} \\
$R'$ & thermal resistance referred to inner area & \si{\meter\squared\kelvin\per\watt} \\
$R''_{f,i}$ & internal fouling resistance & \si{\meter\squared\kelvin\per\watt} \\
$t_{\mathrm{ch}}$, $t_{\mathrm{dc}}$ & charging, discharging duration & \si{\hour} \\
$\Tm$ & PCM melting temperature & \si{\kelvin} \\
$T_{\mathrm{pcm}}$ & PCM temperature, recovered from $E'$ & \si{\kelvin} \\
$\Tre$ & tube outer-surface temperature & \si{\kelvin} \\
$\Ui$ & overall coefficient on inner area & \si{\watt\per\meter\squared\per\kelvin} \\
$V_{\mathrm{well}}$, $\Vmelt$ & PCM volume, melted volume & \si{\meter\cubed} \\
$\delta$ & melt-layer thickness & \si{\meter} \\
$\delta_{\mathrm{merge}}$ & thickness at which neighbouring fronts touch & \si{\meter} \\
$\Delta z$ & segment length & \si{\meter} \\
$\varepsilon_{\mathrm{PCM}}$ & melted fraction of the PCM inventory & -- \\
$\varepsilon_{\mathrm{local}}$ & melted fraction of one segment's PCM & -- \\
$\etaf$, $\etao$ & fin and overall surface efficiency & -- \\
$\lambda$ & assumed storage-loss surplus & -- \\
$\rhom$ & PCM density & \si{\kilogram\per\meter\cubed} \\
\bottomrule
\end{longtable}

Subscripts $s$ and $l$ denote the solid and liquid PCM phases. Subscripts
$\mathrm{ch}$ and $\mathrm{dc}$ denote charging and discharging. A prime on a
resistance denotes referral to the inner tube area, and a prime on a heat rate
denotes per unit tube length.
